\section{The computation model}\label{sec:formalism}

This section fixes the conventions used by the construction and by its
lower bound. All input structures are finite and relational. No order of
their universes is available unless it is explicitly defined from the input.

\subsection{Hereditarily finite sets and ordinary iteration}

For a finite set \(A\) of atoms, \(\HF(A)\) contains the atoms and all
hereditarily finite sets built from them. Atoms are distinct from sets and
have no members. A permutation of \(A\) acts recursively on \(\HF(A)\).
For a set \(a\), let \(\TC(a)\) be the set of all objects reachable from
\(a\) by a positive-length downward membership chain. In particular, the
root \(a\) is not included by a length-zero chain. Put \(\TC(a)=\emptyset\)
for an atom. We use \(1+|\TC(a)|\) as a hereditary-size bound for a set,
and size \(1\) for an atom.

The primitive functions of the BGS syntax are
\[
 \emptyset,\quad \Atoms,\quad \Pair,\quad \Union,\quad \Unique,\quad \Card .
\]
Here \(\Atoms=A\), \(\Pair(a,b)=\{a,b\}\), and
\(\Union(a)=\{b:\exists c\in a,\ b\in c\}\).
The function \(\Unique\) returns the member of a singleton set, and
\(\emptyset\) otherwise. The function \(\Card\) returns the finite von
Neumann ordinal \(|a|\) for a set \(a\), and \(\emptyset\) for an atom.
Input relations are false on tuples having a non-atom entry.

Terms are built from these functions, variables, comprehensions
\[
 \{\,s(\bar a,u):u\in t(\bar a),\ \theta(\bar a,u)\,\},
\]
and iteration from \(\emptyset\). Formulas use input relations, equality
and Boolean combinations; bounded quantifiers are abbreviations using
comprehensions. If the value of a comprehension's range term is an atom,
it has no members and the comprehension is empty.

For an iteration with step \(s(\bar a,x)\), write
\[
 a_0=\emptyset,\qquad a_{i+1}=s(\bar a,a_i).
\]
Its unbounded value is the first stationary value \(a_\ell=a_{\ell+1}\),
or \(\emptyset\) if no such index exists. A fixed polynomial annotation
\(p\) gives the ordinary \(\CPT\) value \(a_\ell\) if
\(\ell\le p(|A|)\) and \(|\TC(a_i)|\le p(|A|)\) for every \(i\ge0\);
otherwise the value is \(\emptyset\). The same bounded semantics is used
recursively in the step term. In particular a rejected candidate and an
empty-set default must be distinguished in a description of evaluation.

We use fixed-length tuple codes formed with ordered pairs, finite products,
set difference, union, bounded minima and case definitions as abbreviations.
For example, a bounded existential formula tests whether its witness
comprehension is nonempty, and a product is obtained from nested
comprehensions followed by finite unions. All abbreviations below admit
direct capture-avoiding substitution into this syntax. The bounds are
verified after choosing such a fixed expansion.

\subsection{Witnessed choice and true-model classes}

We use the WSC operator of~\cite[Sections~3.1--3.2]{ls2023},
written
\[
 \WSC^{*}xy.
 \bigl(s_{\mathrm{step}}(\bar z,x,y),
       s_{\mathrm{choice}}(\bar z,x),
       s_{\mathrm{wit}}(\bar z,x,y),
       \theta_{\mathrm{out}}(\bar z,x)\bigr).
\]
All functions defining these terms are equivariant: applying an input
isomorphism to their arguments applies it to their values.

For fixed parameters \(\bar a\), a choice computation starts at
\(b_0=\emptyset\). From \(b_i\), it chooses an element \(c_i\) of
\(s_{\mathrm{choice}}(\bar a,b_i)\) and sets
\(b_{i+1}=s_{\mathrm{step}}(\bar a,b_i,c_i)\).
Our choice sets will always be nonempty. A first-stabilization path ends
with the first adjacent repetition; if its terminal value is \(b_t\),
the witnessing term is evaluated at
\[
             s_{\mathrm{wit}}(\bar a,b_t,b_i)
\]
for each preceding stage \(i\). Thus the two slots after the parameters
mean \emph{terminal value} and \emph{intermediate value} in the witnessing
term, although they mean current state and choice in the step term.

A set \(M_i\) of permutation graphs witnesses a nonempty choice set \(D_i\)
at stage \(i\) if
\[
 M_i\subseteq\Aut(A,\bar a,b_0,\ldots,b_i)
 \quad\text{and}\quad
 \forall u,v\in D_i\ \exists g\in M_i\ (g(u)=v).
 \tag{*}\label{mod:witness-condition}
\]
The notation on the right of the inclusion means automorphisms of the
\emph{entire relational input} fixing the displayed HF objects.
Permutation graphs act on HF objects by their recursive extensions.
Equivariance makes \(D_i\) invariant under this stabilizer, so
\eqref{mod:witness-condition} makes it an orbit.

Every first-stabilization path must be witnessed. If evaluation succeeds,
the operator returns true when the output formula holds at every terminal
fixed point, and false otherwise. Under successful witnessing these
fixed points lie in one orbit over the parameters, so
their invariant output values agree. We will verify witnessing directly
for every reachable state, rather than depend on a favorable computation
path.

The polynomial annotation bounds the length and the state transitive
sizes of every first-stabilization path. We will also bound every
individual temporary HF object constructed while evaluating its terms.
A failed witness or excessive WSC resource gives \(\dagger\), which
propagates through enclosing expressions. In particular \(\dagger\) is
not changed to false inside a comprehension or Boolean connective.
The collection of all terminal fixed points need not be constructed as
one polynomial-size HF set.

For either formalism \(\mathcal L\), let
\[
 \operatorname{Def}_{\mathcal L}(\tau)
 =\left\{\{A:A\text{ is a finite }\tau\text{-structure and }
                  \varphi(A)=\mathrm{true}\}:
                 \varphi\text{ is an }\mathcal L[\tau]\text{-sentence}\right\}.
\]
This is the comparison in Theorem~\ref{thm:main}. The polynomial, the
vocabulary and the syntax of the separating sentence are fixed before
the input size varies. The full syntax may nest witnessed choices; the
sentence constructed here needs only one occurrence.
